\subsection{Resultados e análise}

As saídas híbridas coincidiram com a saída sequencial em todas as
144 execuções paralelas, incluindo os aquecimentos. Os checksums,
da carga pequena à grande, foram
\texttt{17073017104646724864}, \texttt{17187382406529706254} e
\texttt{14825089465781038537}.

A configuração $2\times4$ em dois nós apresentou a menor média nas
três cargas (Tabelas~\ref{tab:lasdpc-pequena} a~\ref{tab:lasdpc-grande}).
Na carga grande, o tempo médio passou de 13,28~s no sequencial para
2,05~s, com speedup de 6,48 e eficiência de 81,0\%. A diferença entre
as cargas aparece na eficiência dessa configuração, de 55,2\% na
pequena a 81,0\% na grande.

Com dois ranks e duas threads por rank, calculamos a diferença
de tempo entre dois nós e um nó em cada rodada. Na carga média,
a diferença média foi de $-0{,}0228$~s, com IC de 95\% de
$[-0{,}0344; -0{,}0112]$~s. Nas cargas pequena e grande, os
intervalos dessas diferenças incluíram zero. Como a colocação
muda tanto os caminhos de comunicação quanto os recursos que cada
rank compartilha, a comparação diz respeito à configuração completa.

A Figura~\ref{fig:tb2-desempenho} reúne tempo, speedup e eficiência.
Nos rótulos, \texttt{one} identifica o \texttt{hal01} e
\texttt{two} identifica \texttt{hal01} e \texttt{hal02}.
